% =============================================================================
%  1-Introduccion.tex — Introducción / Planteamiento del problema
% =============================================================================

\section{Introducción}
\label{sec:introduccion}

El tema de tesis, aunque suena muy largo y técnico, analiza la viabilidad de los corredores
con geometría anillar. La primera propuesta se basa en el lugar geográfico conocido como
\textit{Anillo Periférico}, que corre dentro de la \zmvm en la Ciudad de México y el Estado
de México. El porqué se eligió esta zona se deriva de los planes de desarrollo que existieron
en 1940 y 1960, en los cuales se definió el concepto de \termino{Los doce Cuarteles}, que
delimitan el centro corporativo y financiero de la capital. Sin embargo, el crecimiento de la
zona en la segunda mitad del siglo XX hizo que el transporte fuera superado por su población.

El enfoque de la tesis es cuantitativo y, a diferencia de otros casos de investigación donde
se usan análisis del ámbito urbanístico como documentos demográficos, encuestas o informes
económicos, se emplea un enfoque matemático-computacional-urbano, fuertemente centrado en el
modelado de entes matemáticos conocidos como \termino{Grafos Matemáticos} para la propuesta
del corredor anillar. Sin embargo, para llegar a dicha propuesta se necesita hacer primero un
análisis de la red de transporte actual (2026), por lo que la investigación se divide en tres
etapas: teoría, introducción y aclaración de conceptos; construcción del modelo matemático basado
en el estado actual de la red; y construcción de la propuesta mediante pruebas de estrés sobre el
modelado de la red actual con el transporte anillar propuesto.

\subsection{Contexto del Semestre 2026-2}
\label{ssec:contexto-semestre}

Durante los primeros días de febrero de 2026 sufrí un problema de salud que me llevó
a estar hospitalizado un par de días en el Instituto de Psiquiatría José Ramón de la Fuente Muñiz,
por lo que me fue imposible acudir presencialmente a las clases del semestre 2026-2
(enero a junio). Aunque solo estuve internado algunos días y en su mayoría fuera de los días
hábiles, la recuperación fue muy lenta y me vi incapacitado para continuar mis actividades durante
los meses de febrero y marzo. Si bien pude comunicarme con la maestra Ángeles Puente e Isabel Reyes,
yo mismo desconocía el diagnóstico y cómo sería el tratamiento de mi enfermedad, lo que hizo que,
aun cuando me sintiera mejor, mi estado de salud permaneciera delicado. Pude retomar mis actividades
académicas en línea a finales de marzo y principios de abril de 2026; por esa razón el avance de
mi investigación quedó rezagado respecto a lo que ya tenía trabajado con anterioridad.